\usepackage{caption}
\usepackage{fancyhdr}
\usepackage{tikz}

\newlength{\runestroke}
\setlength{\runestroke}{0.4pt}

\tikzset{
	rune/.style={
		x=1em,
		y=1em,
		line width=\runestroke
	}
}

% Base outline, always present
\newcommand{\runebase}{%
	\draw (-0.25, 0) -- (0, 0.25) -- (0.25, 0) -- (0, -0.25) -- cycle;
}

% Bit 0, inner diamond
\newcommand{\runefeaturezero}{%
	\fill (-0.245, 0) -- (0, 0.245) -- (0.245, 0) -- (0, -0.245) -- cycle;
}

% Bit 1, outer diamond
\newcommand{\runefeatureone}{%
	\draw (-0.5, 0) -- (0, 0.5) -- (0.5, 0) -- (0, -0.5) -- cycle;
}

% Bit 2, top-left wing
\newcommand{\runefeaturetwo}{%
	\draw (-0.25, 0.0075) -- (-0.25, 0.5) -- (0, 0.75) -- (0, 0.2575) -- cycle;
}

% Bit 3, bottom-right wing
\newcommand{\runefeaturethree}{%
	\draw (0.25, -0.0075) -- (0.25, -0.5) -- (0, -0.75) -- (0, -0.2575) -- cycle;
}

% Bit 4, top-right wing
\newcommand{\runefeaturefour}{%
	\draw (0, 0.25) -- (0, 1) -- (0.5, 0.5) -- (0.5, 0) -- (0.25, 0) -- cycle;
}

% Bit 5, bottom-left wing
\newcommand{\runefeaturefive}{%
	\draw (0, -0.25) -- (0, -1) -- (-0.5, -0.5) -- (-0.5, 0) -- (-0.25, 0) -- cycle;
}

% Draw one feature by its bit index
\newcommand{\runedrawbit}[1]{%
	\ifcase#1\relax
	\runefeaturezero
	\or
	\runefeatureone
	\or
	\runefeaturetwo
	\or
	\runefeaturethree
	\or
	\runefeaturefour
	\or
	\runefeaturefive
	\fi
}

% Extra shapes for combinations
\newcommand{\runedrawcombinations}{%
	\ifnum\runeflag{1}=1\relax
	\ifnum\runeflag{2}=1\relax
	\fill (-0.25, 0.0075) -- (-0.25, 0.25) -- (0, 0.5) -- (0, 0.2575) -- cycle;
	\fi
	\ifnum\runeflag{3}=1\relax
	\fill (0.25, -0.0075) -- (0.25, -0.25) -- (0, -0.5) -- (0, -0.2575) -- cycle;
	\fi
	\fi
}

% Read whether a bit was selected
\newcommand{\runeflag}[1]{%
	\csname runeFlag#1\endcsname
}

% Parsing and integer conversion
\ExplSyntaxOn

\seq_new:N \l__rune_digits_seq

\seq_new:N \l__rune_bits_seq
\tl_new:N  \l__rune_bits_tl
\int_new:N \l__rune_value_int
\int_new:N \l__rune_index_int

% Explicit space-separated bit indices
\NewDocumentCommand{\runebit}{m}
{
	\group_begin:
	\int_step_inline:nnn {0} {5}
	{
		\cs_set:cpn {runeFlag##1} {0}
	}
	
	\seq_set_split:Nnn \l__rune_bits_seq { ~ } {#1}
	\seq_remove_all:Nn \l__rune_bits_seq {}
	
	\seq_map_inline:Nn \l__rune_bits_seq
	{
		\cs_set:cpn {runeFlag##1} {1}
	}
	
	\tikz[rune, x=1em, y=1em, baseline=-0.35em]
	{
		\runebase
		\seq_map_inline:Nn \l__rune_bits_seq
		{
			\runedrawbit{##1}
		}
		\runedrawcombinations
	}
	\group_end:
}

\DeclareRobustCommand{\runeglyph}[1]
{
	\group_begin:
	\int_set:Nn \l__rune_value_int {#1}
	\int_zero:N \l__rune_index_int
	\tl_clear:N \l__rune_bits_tl
	
	\int_while_do:nn { \l__rune_value_int > 0 }
	{
		\int_if_odd:nT {\l__rune_value_int}
		{
			\tl_put_right:Nx \l__rune_bits_tl
			{
				\int_use:N \l__rune_index_int
				\c_space_tl
			}
		}
		
		\int_set:Nn \l__rune_value_int
		{
			\int_div_truncate:nn
			{\l__rune_value_int} {2}
		}
		\int_incr:N \l__rune_index_int
	}
	
	\exp_args:NV \runebit \l__rune_bits_tl
	\group_end:
}

\DeclareRobustCommand{\runeadd}{%
	\mathord{%
		\tikz[rune, x=1em, y=1em, baseline=-0.35em]{%
			\draw (-0.3, 0) .. controls (0.05, 0.45) and (-0.05, -0.45) .. (0.3, 0);
			\draw (0, 0) circle[radius=0.175];
		}%
	}%
}

\DeclareRobustCommand{\runesub}{%
	\mathord{%
		\tikz[rune, x=1em, y=1em, baseline=-0.35em]{%
			\draw (-0.3, 0) -- (0.3, 0);
		}%
	}%
}

\tl_new:N \l__rune_decimal_tl
\tl_new:N \l__rune_quotient_tl
\int_new:N \l__rune_remainder_int
\int_new:N \l__rune_quotient_digit_int

\DeclareRobustCommand{\rune}[1]
{
	\group_begin:
	\tl_set:Nx \l__rune_decimal_tl {#1}
	\tl_trim_spaces:N \l__rune_decimal_tl
	\seq_clear:N \l__rune_digits_seq
	
	\bool_do_until:nn
	{
		\str_if_eq_p:ee
		{\l__rune_decimal_tl} {0}
	}
	{
		% Divide the decimal string by 64.
		\tl_clear:N \l__rune_quotient_tl
		\int_zero:N \l__rune_remainder_int
		
		\tl_map_inline:Nn \l__rune_decimal_tl
		{
			\int_set:Nn \l__rune_remainder_int
			{
				10 * \l__rune_remainder_int + ##1
			}
			
			\int_set:Nn \l__rune_quotient_digit_int
			{
				\int_div_truncate:nn
				{\l__rune_remainder_int} {64}
			}
			
			% Suppress leading zeros in the quotient.
			\int_compare:nNnTF
			{\l__rune_quotient_digit_int} > {0}
			{
				\tl_put_right:Nx \l__rune_quotient_tl
				{
					\int_use:N
					\l__rune_quotient_digit_int
				}
			}
			{
				\tl_if_empty:NF \l__rune_quotient_tl
				{
					\tl_put_right:Nn
					\l__rune_quotient_tl {0}
				}
			}
			
			\int_set:Nn \l__rune_remainder_int
			{
				\int_mod:nn
				{\l__rune_remainder_int} {64}
			}
		}
		
		% Remainder is the next base-64 digit.
		\seq_put_right:Nx \l__rune_digits_seq
		{
			\int_use:N \l__rune_remainder_int
		}
		
		\tl_if_empty:NTF \l__rune_quotient_tl
		{
			\tl_set:Nn \l__rune_decimal_tl {0}
		}
		{
			\tl_set_eq:NN
			\l__rune_decimal_tl
			\l__rune_quotient_tl
		}
	}
	
	\seq_map_indexed_inline:Nn \l__rune_digits_seq
	{
		\int_compare:nNnT {##1} > {1}
		{
			\kern-\runestroke
		}
		\runeglyph{##2}
	}
	\group_end:
}

\ExplSyntaxOff

\setcounter{section}{-1}
\renewcommand{\thesection}{%
	\protect\rune{\arabic{section}}%
}

\setcounter{table}{-1}
\renewcommand{\thetable}{%
	\protect\rune{\arabic{table}}%
}

\setcounter{page}{0}
\fancypagestyle{plain}{%
	\fancyhf{}
	\fancyhfoffset{0pt}
	\fancyfoot[C]{\rune{\arabic{page}}}
	\renewcommand{\headrulewidth}{0pt}
	\renewcommand{\footrulewidth}{0pt}
}

\pagestyle{plain}